% 卒業研究 発表レジュメ
% コンパイル: lualatex resume.tex または xelatex resume.tex
% （Overleaf ではコンパイラを LuaLaTeX か XeLaTeX に設定）
\documentclass[a4paper,10pt,twocolumn,autodetect-engine,ja=standard]{bxjsarticle}

\geometry{top=15mm,bottom=15mm,left=16mm,right=16mm,columnsep=7mm}
\usepackage{amsmath,amssymb,bm}
\usepackage{booktabs}
\usepackage{enumitem}
\usepackage{tikz}
\usetikzlibrary{positioning,arrows.meta,calc}
\usepackage{titlesec}
\usepackage{indentfirst}

% 節見出しをテンプレートに合わせて小さめ・太字に
\titleformat{\section}{\normalsize\sffamily\bfseries}{\thesection}{1em}{}
\titlespacing*{\section}{0pt}{1.0ex}{0.4ex}
\setlist{nosep,leftmargin=1.2em}
\renewcommand{\labelitemi}{\textbullet}
\setlength{\abovedisplayskip}{3pt}
\setlength{\belowdisplayskip}{3pt}
\pagestyle{plain}
\setlength{\textfloatsep}{8pt plus 2pt minus 2pt}
\setlength{\abovecaptionskip}{3pt}
\setlength{\belowcaptionskip}{0pt}

\newcommand{\x}{\bm{x}}
\newcommand{\thetab}{\bm{\theta}}

\begin{document}

\twocolumn[
  \begin{center}
    \sffamily\bfseries
    {\large 中小企業の製造現場における
     タスク自動割当システムと、実績から学習する割当 AI の開発\par}
    
    \normalfont
    % ↓ 氏名・指導教員名を記入
    E2209 小笠原​悠太\qquad 指導教員：野村隼人
  \end{center}
  \vspace{1mm}
]

\section{はじめに}
工場の共有機材（レーザー加工機・旋盤・プレス機など）では、「誰が何に使っているか
分からない」「予約はあるのに来ない（ゴースト予約）」「重要な仕事が特定の人に偏る・
忘れられる」という非効率が起きやすい。製造業では DX を進める人材の不足が
課題とされており\,\cite{monozukuri,chusho}、専任の IT 担当者がいない中小工場では、
計画を立てるコストと、空いている機材や仕事を探すコストが現場の負担になる。

本研究では、各機材に NFC リーダー付きのモジュールを置き、社員証をタッチした瞬間に
\textbf{その人が今その場所でやるべきタスク}を提示するシステムを開発した。
あえて遠隔予約の機能を持たず、物理的なタッチだけを操作の起点にすることで、
システム上の状態と現場の実態が常に一致する（ゴースト予約が原理的に起きない）。

チームでの開発のうち、筆者は\textbf{タスク割当 AI（WariAthena）}を担当した。
割当は「誰にどの仕事を任せると速く終わるか」を、作業ログから学習しながら決める
必要がある。本稿の貢献は次の 3 点である。
\begin{itemize}
  \item \textbf{システム}（2 節）：NFC タッチ起点の割当・排他制御・誘導
  \item \textbf{割当 AI}（3 節）：作業者ごとのベイズ線形回帰と Thompson Sampling\,\cite{honda,agrawal} による
        文脈付きバンディットとして割当を定式化し、手持ちの仕事量による負荷分散と、
        資格によるハード制約を組み込んだ
  \item \textbf{評価}（4 節）：シミュレーションで経験年数順などの割当と比べた
\end{itemize}

\section{システム全体の構成}
\begin{figure}[t]
  \centering
  \begin{tikzpicture}[
      font=\scriptsize,
      box/.style={draw,rounded corners=2pt,align=center,inner sep=3pt,minimum height=7mm},
      arr/.style={-{Stealth[length=4pt]},semithick}]
    \node[box] (mod) {機材モジュール ×台数\\NFC・LCD・ボタン・LED};
    \node[box,right=6mm of mod] (srv) {サーバ\\Flask + SQLite};
    \node[box,right=6mm of srv] (web) {管理画面\\（ブラウザ）};
    \node[box,below=3mm of srv,fill=gray!12] (ai) {\textbf{割当 AI（WariAthena）}};
    \draw[arr,<->] (mod) -- (srv);
    \draw[arr,<->] (srv) -- (web);
    \draw[arr,<->] (srv) -- (ai);
  \end{tikzpicture}
  \caption{システム構成（モジュール–サーバ間は MQTT）}
  \label{fig:system}
\end{figure}

\textbf{機材モジュール}（図\ref{fig:system}）：試作機は Raspberry Pi 4 に NFC リーダー（RC522）、
2.8 インチ LCD（ILI9341）、物理ボタン 3 個（左・決定・右）、状態 LED（作業中は赤、
それ以外は青）をつないだ。量産は ESP32 へ移す計画である。
\textbf{通信}：サーバとは MQTT\,\cite{mqtt} で双方向に通信する。ブローカーは自動探索し、
切断は Last Will で即時に死活表示へ反映する。
\textbf{サーバ}：Flask と SQLite。

\textbf{現場の流れ}：(1) 空いている機材で社員証をタッチすると、候補タスクが 1 件ずつ
LCD に出る。(2) ボタンで承認すると機材がロックされ、他者のタッチは拒否される（排他制御）。
候補が無い・すべて断ったときは「フリー利用」としてロックする。(3) 他の機材に本人宛の
より重要なタスクがあれば、そちらへ移るかを尋ねる（誘導）。(4) 再タッチで完了し、
所要時間を作業ログに記録して機材を解放する。直後に体感難易度（簡単・普通・難しい）を尋ねる。

\textbf{管理画面}：タスクの登録（重要度・難易度・数量・期限・必要な資格）、
AI による割当とその根拠の表示、稼働状況のダッシュボードを持つ。
画面からの機材予約は、ゴースト予約を防ぐため意図的に作っていない。

\section{割当 AI：WariAthena}
\textbf{問題設定}：割当は、作業者を「腕」、タスクの特徴を「文脈」とする
文脈付きバンディット\,\cite{honda} とみなせる。任せたことのない人の実力は分からないため、実績のある人に任せる\textbf{活用}と、
まだ分からない人に任せてみる\textbf{探索}の両立が要る。

\textbf{文脈}：タスク $j$ を作業者 $w$ に任せるときの文脈ベクトル $\x$ を
\begin{equation}
  \x = \bigl[\,y,\; d,\; y d,\; \bm{e}_{\mathrm{diff}},\; \bm{e}_{\mathrm{perm}},\; \bm{e}_{\mathrm{eq}}\,\bigr]
\end{equation}
とした。$y$ は勤続年数 / 30、$d$ は難易度（1〜5 を 0.15〜0.85 に写す）、
$\bm{e}_{\mathrm{diff}}$ は難易度の one-hot、$\bm{e}_{\mathrm{perm}}$ は必要な資格
（溶接・検査など）の multi-hot、$\bm{e}_{\mathrm{eq}}$ は機材の one-hot である。
資格と機材を入れたのは、難易度だけでは「溶接が得意な人」「旋盤が得意な人」を区別できないためである。

\textbf{報酬}：作業ログから、1 個あたりの所要時間 $t$（数量で割る）と、
同じ難易度の過去の平均 $\bar t$ を用いて
\begin{equation}
  r = c_f \cdot \bar t / t, \qquad c_f \in \{1.5,\ 1.0,\ 0.5\}
\end{equation}
とした。$c_f$ は本人の体感難易度（簡単・普通・難しい）の係数で、未回答なら
$t$ が $\bar t$ より 1 割以上速い・遅いかで決める。数量で割らないと、
大量のタスクを引いた人が「遅い」と学習されてしまう。

\textbf{学習}：作業者ごとに $r \approx \thetab_w^{\top}\x$ の線形モデルを置き、
事前分布 $\thetab_w \sim \mathcal N(\bm 0, 100 I)$、観測雑音 $\sigma^2 = 0.1$ の
ベイズ線形回帰\,\cite{suyama,bishop} で事後分布 $\mathcal N(\bm\mu_w, \Sigma_w)$ を更新する：
\begin{align}
  \Sigma_w^{-1} &\leftarrow \Sigma_w^{-1} + \sigma^{-2}\x\x^{\top}, \\
  \bm\mu_w &\leftarrow \Sigma_w\bigl(\Sigma_{w,\mathrm{old}}^{-1}\bm\mu_{w,\mathrm{old}} + \sigma^{-2} r\,\x\bigr).
\end{align}
更新は足し算なので、学習結果は保存せず、作業ログを古い順に再生して
組み立て直す（再起動しても失われない）。

\textbf{選択（Thompson Sampling）}：各作業者の事後分布から $\tilde\thetab_w$ を 1 本ずつ引き、
予測報酬 $\tilde r_w = \tilde\thetab_w^{\top}\x$ で比べる\,\cite{honda,agrawal}。実績の少ない人は
分布が広く、ときどき大きな値を引くので自然に試される。

\textbf{負荷分散}：予測報酬だけで選ぶと、何でも速いベテランに仕事が集まる。
そこで作業者 $w$ の手持ちの残り時間 $L_w$ と、このタスクの見積もり時間 $\hat T_j$
（難易度ごとの平均 × 数量）を用いて
\begin{equation}
  s_w = \tilde r_w \Big/ \bigl(1 + L_w / \hat T_j\bigr)
\end{equation}
で選ぶ（$\tilde r_w < 0$ のときは掛ける）。引き算でなく割り算にしたのは、
$c_f$ のため報酬が速さに比例せず尺度が一定でないからで、割り算なら尺度に依らず、
得意・不得意の差が大きい仕事ほど任され続ける。

\textbf{資格はハード制約}：資格（フォークリフト・クレーン・溶接など）は学習させず、
AI を呼ぶ前に候補から除く。無資格者に危険作業を任せて失敗から学ぶ設計は成り立たず、
勤続年数でも代用できない。

\textbf{使い方は 2 通り}：(a) 現場でタッチされたとき、その人の $\tilde\thetab_w$ を
1 本だけ引き、候補タスクを予測報酬の順に並べる。
(b) 管理画面の「AI 割当」では、タスクに対し上の $s_w$ が最大の人を選び、
各候補の期待値・引いた値・手持ち・除外理由を根拠として画面に残す。

\section{評価}
\textbf{設定}：作業者 6 名（勤続 1〜30 年）に、真の時間倍率を隠して与えた。
ベテランほど全体に速いが（倍率 $1.3 - 0.5\,y$）、個人差 $\pm 30\%$ があり、
各人に 1 つ得意な作業（溶接・検査・プレスのいずれか、倍率 $\times 0.6$）がある。
タスクは難易度 1〜5（標準 1〜10 分/個）、数量 1〜5、種類 4 通りを一様に生成し、
1 ラウンド 6 件を順に割り当て、ラウンドの終わりに全員が作業を終えてログに加わる。
150 ラウンド、20 試行の平均をとった。割当 AI は実装をそのまま呼んだ。

\textbf{指標}：(i) \textbf{適性比}：任せた人の時間倍率 / そのタスクに最適な人の倍率
（1 が最良）、(ii) \textbf{完了時間}：1 ラウンドで最も遅い人が終わるまでの時間、
(iii) \textbf{偏り}：最も多く担当した人の割合。(i)(ii) は学習後（51 ラウンド目以降）で測った。

\begin{table}[t]
  \centering
  \caption{割当方式の比較（20 試行の平均）}
  \label{tab:result}
  \small
  \begin{tabular}{lccc}
    \toprule
    方式 & 適性比 & 完了時間 [分] & 偏り \\
    \midrule
    ランダム           & 1.73 & 44.4 & 19\% \\
    順番（ラウンドロビン） & 1.72 & 34.7 & 17\% \\
    経験年数が最長の人     & 1.27 & 63.7 & 100\% \\
    提案（負荷分散なし）   & \textbf{1.03} & 32.4 & 42\% \\
    \textbf{提案}         & 1.16 & \textbf{25.4} & 27\% \\
    \midrule
    参考：真の適性を知る   & 1.21 & 23.6 & 24\% \\
    \bottomrule
  \end{tabular}\\[2pt]
  {\footnotesize 参考：真の倍率を知り、最も早く終わる人へ任せる貪欲法}
\end{table}

\textbf{結果}（表\ref{tab:result}）：経験年数で選ぶと全件がベテラン 1 人に集まり
（偏り 100\%）、他の人が遊ぶため完了時間が最も長い。ランダムや順番は得意でない人に任せるので適性比が 1.7 と悪い。

負荷分散なしの提案は、学習後はほぼ常に最適な人を選べる（適性比 1.03）。
作業ログだけから隠れた得意分野を見つけられている。
ただし得意な人に仕事が集まり（偏り 42\%）、完了時間は順番とあまり変わらない。

負荷分散を加えた提案は、得意な人が埋まれば次に向く人へ回すため適性比は 1.16 に上がるが、完了時間は 25.4 分で、経験年数による割当より 60\%、
順番より 27\% 短い。真の適性を知る貪欲法（23.6 分）との差は 8\% である。
「各人に一番向く仕事」と「全体が早く終わる割当」は一致せず、後者には負荷分散が要る。

\section{まとめと今後の展望}
NFC タッチを起点とするタスク割当システムと、作業ログから各人の得意・不得意を学び、
負荷を分散しながら割り当てる AI を開発した。
シミュレーションでは、真の適性を知る割当に近い完了時間を得た。

\textbf{限界}：評価はシミュレーションのみで、実際の工場の作業ログはまだ無い。
作業者の適性は時間とともに変わらないと仮定した（実際は熟練で速くなる）。
ログから毎回学習し直すので、ログが増えると計算が重くなる。
「指定した人に任せる」タスクは未対応である。

\textbf{今後の展望}：実際の現場でのログの収集と評価、熟練による変化を扱うための
古いログの割引、役職の特徴量化、ESP32 版モジュールの製作を進める。

\footnotesize
\begin{thebibliography}{9}\setlength{\itemsep}{0pt}
  \bibitem{monozukuri} 経済産業省 他，『2025 年版ものづくり白書』，2025．
  \bibitem{chusho} 中小企業庁，『2025 年版中小企業白書』，2025．
  \bibitem{honda} 本多淳也・中村篤祥，『バンディット問題の理論とアルゴリズム』，講談社，2016．
  \bibitem{suyama} 須山敦志，『ベイズ推論による機械学習入門』，講談社，2017．
  \bibitem{bishop} C. M. ビショップ，『パターン認識と機械学習（上）』，丸善出版，2012．
  \bibitem{agrawal} S. Agrawal and N. Goyal, ``Thompson sampling for contextual bandits with linear payoffs,'' \textit{Proc. ICML}, 2013.
  \bibitem{mqtt} OASIS, ``MQTT Version 5.0,'' OASIS Standard, 2019.
\end{thebibliography}

\end{document}
